\documentclass{article}

\usepackage[T1]{fontenc}
\usepackage[brazilian]{babel}
\usepackage{hyperref}
\usepackage[
  profundidade-secoes=5,
  profundidade-sumario=3
]{abntex3}

\ABNTEXstructureSetup{1}{destaque=negrito}
\ABNTEXstructureSetup{2}{destaque=negrito}
\ABNTEXstructureSetup{3}{destaque=italico}

\begin{document}

\ABNTEXpretextualtitle{Agradecimentos}

Este exemplo agradece às pessoas que colaboram com o desenvolvimento aberto.

\ABNTEXsumario

\section{Introdução à estrutura}

O primeiro nível inicia a numeração progressiva.

\subsection{Uma seção com título suficientemente longo para demonstrar o
alinhamento das linhas seguintes}

O texto começa em uma nova linha e permanece relacionado ao título.

\subsubsection{Ações, órgãos e informações}

Os caracteres acentuados também exercitam os marcadores de PDF quando
\texttt{hyperref} está carregado.

\paragraph{Quarto nível}

O quarto nível continua a hierarquia.

\subparagraph{Quinto nível}

O quinto nível encerra a profundidade permitida.

\ABNTEXposttextualtitle{Referências}

As referências serão integradas ao \texttt{biblatex-abnt} em marco posterior.

\ABNTEXappendix{Questionário de avaliação}

Material elaborado pelo autor.

\ABNTEXannex{Autorização institucional}

Material de autoria externa incorporado ao trabalho.

\end{document}
